\section{Weighted coefficients and actual expansions}
\label{ads:ana:section}

We first justify the passage between finite Taylor calculations and the
given solution branch.  Throughout this section, fix
\begin{equation}
  \frac32<\beta<\min\{\tau,\sqrt3\}.
  \label{ads:ana:beta}
\end{equation}
We use the weighted norms and Taylor conventions of
Section~\ref{ads:sec:taylor}. Boundary norms without a weight are uniform
norms on $S$.

\subsection{A scalar inverse and its asymptotics}

Equivariant linear differential operations preserve finite angular type.  Products and contractions preserve
it after enlarging the angular space by a finite tensor product.  On
functions, the round Laplacian preserves these spaces: it is the sum of
squares of the three infinitesimal rotation fields.  Its restriction is
self-adjoint, so we may split into finitely many eigenspaces with
eigenvalues $\lambda\geq0$ for $-\Delta_\sigma$.

\Needspace{9\baselineskip}
\begin{lemma}[Comparison and finite angular inverse]
\label{ads:lem:inverse}
There are constants $c_\beta,C_\beta>0$ with the following properties.
Suppose $z=o(e^{-\beta s})$ uniformly at infinity and
\[
  (L-V)z=F,\qquad \|V\|_{C^0(M)}\leq c_\beta.
\]
If either $z|_S=b$ or $\partial_s z|_S=b$, then
\begin{equation}
  \|z\|_\beta\leq C_\beta
       \bigl(\|F\|_\beta+\|b\|_{C^0(S)}\bigr).
  \label{ads:ana:comparison}
\end{equation}
The constants do not depend on an angular cutoff.  The same estimate
holds for $\partial_s z-a_*z=b$ if
$\|a_*\|_{C^0(S)}\leq\beta/2$.

If $F$ and $b$ are smooth and of finite angular type, and
$F=O(e^{-\gamma s})$ for some $\gamma>3$, there is a unique decaying
finite angular type solution of $Lu=F$ with either of the above
Dirichlet or Neumann data.  It satisfies
\begin{equation}
  u=r^{-3}\zeta(x)+o(r^{-3}),\qquad
  \partial_s\bigl(u-r^{-3}\zeta(x)\bigr)=o(r^{-3})
  \label{ads:ana:inverse-asymptotic}
\end{equation}
for a smooth finite angular type function $\zeta$.  These statements
hold after any fixed number of angular differentiations.  If
$\partial_s^jF=O(e^{-\gamma s})$ for $0\leq j\leq k$, the remainder
in \eqref{ads:ana:inverse-asymptotic} can also be differentiated in $s$
through order $k+2$, with each resulting remainder $o(r^{-3})$.
\end{lemma}

\begin{proof}
Let $h=e^{-\beta s}$.  Since $N/r\geq0$,
\[
  Lh=(\beta^2-2\beta N/r-3)h
       \leq-(3-\beta^2)h.
\]
Choose $c_\beta<(3-\beta^2)/2$.  A sufficiently large multiple
$Ah$ satisfies $(L-V)(Ah)\leq-|F|$.  For Dirichlet data, enlarge
$A$ to dominate $|b|$.  For Neumann data, enlarge it so that
$-\beta A-b<0$ and $-\beta A+b<0$ on $S$.
The functions $Ah-z$ and $Ah+z$ are positive sufficiently far out.
A negative interior minimum contradicts their differential inequalities,
because the zeroth-order coefficient of $L-V$ is negative.
A negative minimum on $S$ is impossible in the Neumann case: the
inward derivative at such a minimum is nonnegative, whereas the
chosen derivative is negative.  This proves \eqref{ads:ana:comparison}.
For the Robin condition, apply the maximum argument to $z/h$.
At a positive boundary maximum its inward derivative is nonpositive,
whereas
$\partial_s(z/h)=(\beta+a_*)(z/h)+b$ on $S$.
Since $\beta+a_*\geq\beta/2$, a sufficiently large positive maximum
is impossible; apply the same argument to $-z/h$.

For the existence assertion, work in each angular eigenspace.  The
radial equation is
\begin{equation}
 y''+2\frac Nr y'-\left(3+\frac\lambda{r^2}\right)y=F_\lambda.
 \label{ads:ana:radial-ode}
\end{equation}
First impose the required condition at $0$ and $y(T)=0$ on a finite
interval.  The homogeneous problem has trivial kernel: multiplication
by $r^2y$ and integration gives
\[
  0=-\int_0^T r^2(y')^2\,\dd s
    -\int_0^T(3r^2+\lambda)y^2\,\dd s,
\]
with zero boundary terms.  Thus the finite-interval linear problem is
solvable.  The barrier estimate is independent of $T$.  On each compact
interval, the equation bounds $y'$ and then $y''$: the mean value
theorem supplies a point with bounded $y'$, and the first-order equation
for $y'$ propagates that bound.  A subsequence therefore converges on
compact intervals to a solution with $y=O(e^{-\beta s})$.
The same local argument on unit intervals gives
$y'=O(e^{-\beta s})$.

Now $N/r=1+O(e^{-2s})$ and $r^{-2}=O(e^{-2s})$.  Hence
\[
  y''+2y'-3y=G,\qquad
  G=O(e^{-\gamma' s}),\qquad
  \gamma'=\min\{\gamma,\beta+2\}>3.
\]
Variation of constants for the roots $1,-3$ eliminates the growing
homogeneous term and gives
\[
 y=c e^{-3s}+O(e^{-\gamma' s}),\qquad
 y'=-3c e^{-3s}+O(e^{-\gamma' s}).
\]
For example the particular solution is a linear combination of
\[
 e^s\int_s^\infty e^{-t}G(t)\,\dd t
 \quad\text{and}\quad
 e^{-3s}\int_s^\infty e^{3t}G(t)\,\dd t;
\]
both integrals converge
with the claimed bounds.  Since
$r=C e^s(1+O(e^{-2s}))$, this is
\eqref{ads:ana:inverse-asymptotic}.  The equation gives the second
differentiated remainder, and differentiated equations give the rest.
Angular differentiations are bounded operations on the fixed finite
spaces.  Any other solution tending uniformly to zero also satisfies
the initial weighted bound: compare on sufficiently long intervals
with $Ah+\delta$, where $\delta>0$ dominates the far boundary value,
and then let $\delta\downarrow0$.  The same ODE improvement applies.
The improved decay permits comparison for a difference of two
solutions, proving uniqueness.
\end{proof}

\subsection{Coefficient equations and the required spatial derivatives}

Put $\psi=\phi-1$, $h_q=q_{ss}|_S$, and $Q_q=|q|_{g_0}^2$.  Define
\begin{align}
 \mathcal A(z)&=\frac{15}{2}z^2+\frac{15}{2}z^3
                    +\frac{15}{4}z^4+\frac34z^5,\notag\\
 \mathcal N_{\eps,q}(z)&=\mathcal A(z)
                    -\frac{\eps^2}{8}Q_q(1+z)^{-7},\notag\\
 \mathcal B_{\eps,q}(z)&=\frac{\eps}{4}h_q(1+z)^{-3}.
 \label{ads:ana:nonlinearities}
\end{align}
The actual equations become
\begin{equation}
 L\psi=\mathcal N_{\eps,q}(\psi),\qquad
 \partial_s\psi|_S=\mathcal B_{\eps,q}(\psi|_S).
 \label{ads:ana:actual-equations}
\end{equation}

For a finite angular type seed, define $u_j=u_j[q]$ recursively.
Set $P_{j-1}=\sum_{i<j}\eps^iu_i$, and let
\begin{equation}
 \begin{gathered}
 F_j=[\mathcal N_{\eps,q}(P_{j-1})]_j,
 \qquad B_j=[\mathcal B_{\eps,q}(P_{j-1}|_S)]_j,\\
 Lu_j=F_j,\qquad \partial_su_j|_S=B_j,
 \qquad u_j\longrightarrow0.
 \end{gathered}
 \label{ads:ana:coefficient-recursion}
\end{equation}
Here $[\cdot]_j$ extracts the coefficient of $\eps^j$; inverse powers
are Taylor expanded at $1$.  The first two equations are explicitly
\begin{align}
 F_1&=0,& B_1&=\tfrac14h_q,\notag\\
 F_2&=\tfrac{15}{2}u_1^2-\tfrac18Q_q,
 & B_2&=-\tfrac34h_qu_1|_S.
 \label{ads:ana:first-two}
\end{align}
Only preceding coefficients occur on the right of
\eqref{ads:ana:coefficient-recursion}.  The construction does not require
an actual solution branch for the seed.

\begin{lemma}[Finite coefficient asymptotics]
\label{ads:ana:coefficients}
For every smooth finite angular type seed in $C^{1,\alpha}_\tau$,
the coefficients through degree two are well defined and
\[
  u_j=r^{-3}\zeta_j(x)+o(r^{-3}),\qquad j=1,2.
\]
The first coefficient has this asymptotic with arbitrarily many
spatial derivatives.  The second has it with the first spatial
derivatives, in particular with every derivative needed at degree
two below.

For the fourth-order construction, suppose specifically that
\begin{equation}
 \begin{gathered}
 q=\Hess v-Bv,\qquad Lv=0,\qquad v\longrightarrow0,\\
 v|_S\text{ has only constant and degree-one spherical modes}.
 \end{gathered}
 \label{ads:ana:kernel-seed}
\end{equation}
Corollary~\ref{ads:cor:kernel} will show that every seed with zero quadratic
deficit has this form, so fourth-order estimates are needed only in
this class. Then $q$ and all its derivatives have $O(r^{-3})$ scaled
components.
The coefficients $u_1,\ldots,u_4$ exist and have the above asymptotic
with arbitrarily many spatial derivatives.

Let $p=2$ in the first case and $p=4$ in the second.  In the coordinate
$\eta=\operatorname{arsinh}r$, the conformal metric jet
$G=(1+\sum_{j=1}^p\eps^ju_j)^4g_0$, taken through degree $p$, satisfies
\begin{equation}
 \begin{gathered}
 G-b=r^{-3}\bigl(C\,\dd\eta^2+E r^2\sigma\bigr)+o(r^{-3}),\\
 C=2m_0+4\sum_{j=1}^p\eps^j\zeta_j,\qquad
 E=4\sum_{j=1}^p\eps^j\zeta_j.
 \end{gathered}
 \label{ads:ana:metric-asymptotic}
\end{equation}
For coefficient $j$, the remainder has this decay after angular and
$\eta$ derivatives through order $p-j+1$, in scaled coframes
$\dd\eta,r\,\dd x$.
\end{lemma}

\begin{proof}
Lemma~\ref{ads:lem:inverse} constructs $u_1$ with homogeneous interior
equation.  Its derivatives have the asserted asymptotics by
\eqref{ads:ana:radial-ode}.  In degree two the interior source is a sum
of $u_1^2=O(r^{-6})$ and $|q|^2=O(r^{-2\tau})$, with
$\min\{6,2\tau\}>3$.  Its first radial derivative has the same
bound because $q\in C^{1,\alpha}_\tau$.  Finite angular type supplies
any fixed number of angular derivatives.  The inverse lemma therefore
gives the claimed second coefficient and, in particular, its first
differentiated asymptotic.  No bound on higher radial derivatives of a
general seed is asserted here.

For \eqref{ads:ana:kernel-seed}, solve the homogeneous Dirichlet problem
separately in its four boundary modes.  The homogeneous ODE gives
$v=O(r^{-3})$ with all differentiated versions.  The scaled components
of the connection, $B$, and all their derivatives are bounded at
infinity.  It follows that $\Hess v-Bv$ has $O(r^{-3})$ components
with all derivatives.  Inductively, every interior source in
\eqref{ads:ana:coefficient-recursion} is a sum of products containing at
least two factors among lower solution coefficients and $q$.
It therefore has $O(r^{-6})$ decay with all derivatives.  Another
application of the inverse lemma completes the induction through
degree four.

Finally $g_0-b=2m_0r^{-3}\dd\eta^2+O(r^{-5})$ in scaled components,
with all derivatives.  In each positive degree the only $r^{-3}$
term in $(1+\sum\eps^ju_j)^4-1$ is $4r^{-3}\zeta_j$;
products of two solution coefficients are $O(r^{-6})$.
This proves \eqref{ads:ana:metric-asymptotic}.  At $p=2$, the derivative
count is two for coefficient one and one for coefficient two, both
already established.  At $p=4$ all the needed derivatives are
available from the stronger case.
\end{proof}

\subsection{An exact Green formula for the mass}

Define the test functions
\begin{equation}
 W_0=N,\qquad W_i=rx_i\quad(1\leq i\leq3).
 \label{ads:ana:tests}
\end{equation}
The background identities and the radial equation give
\begin{equation}
 LW_0=0,\qquad LW_i=-3m_0r^{-2}x_i.
 \label{ads:ana:tests-equations}
\end{equation}

\begin{lemma}[Green representation of the mass]
\label{ads:lem:green}
For every sufficiently small member of the actual branch,
\begin{equation}
 p_\mu(\eps)-p_\mu(0)
 =\frac1{2\pi}\left\{
   \int_S\bigl(\psi\partial_sW_\mu-W_\mu\partial_s\psi\bigr)\,\dd A_{g_0}
   -\int_M\bigl(W_\mu L\psi-\psi LW_\mu\bigr)\,\dd V_{g_0}
                 \right\}.
 \label{ads:ana:green}
\end{equation}
In particular the right-hand side is an absolutely convergent integral
expression.  More generally the linear functional
\begin{equation}
 \mathcal G_\mu(u,F,B)
 =\frac1{2\pi}\left\{
 \int_S(u\partial_sW_\mu-W_\mu B)\,\dd A_{g_0}
 -\int_M(W_\mu F-uLW_\mu)\,\dd V_{g_0}\right\}
 \label{ads:ana:green-functional}
\end{equation}
is bounded in the norms
$\|u\|_\beta+\|F\|_{2\beta}+\|B\|_{C^0(S)}$.
For finite angular type coefficient jets, this same functional computes
the coefficient of their metric mass.
\end{lemma}

\begin{proof}
First fix $\eps$; no parameter limit is taken in this argument.
The assumed weighted regularity gives
$\psi,\nabla^b\psi=O(r^{-\tau})$.  The difference between the actual
metric perturbation $\phi^4g_0-g_0$ and $4\psi b$ consists of terms
with scaled components and first derivatives
\[
 O(r^{-2\tau})+O(r^{-\tau-3}).
\]
Indeed $\phi^4-1-4\psi=O(\psi^2)$ and
$g_0-b=O_1(r^{-3})$.  Since a mass test function and its derivative
have size $O(r)$ and the sphere area has size $O(r^2)$, these terms
give flux errors
\begin{equation}
 O(r^{3-2\tau})+O(r^{-\tau})\longrightarrow0.
 \label{ads:ana:flux-errors}
\end{equation}
For the tensor $4\psi b$ in dimension three,
\[
 \Div_b(4\psi b)-\dd\tr_b(4\psi b)=-8\dd\psi,
\]
and the remaining two terms of the prescribed mass integrand add
$8\psi\,\dd V_\mu$.  Thus the flux difference is $1/(2\pi)$ times
the flux of $\psi\grad_bV_\mu-V_\mu\grad_b\psi$.
Replacing the hyperbolic normal by $\partial_s$ produces a vanishing
error since the relative normal change is $O(r^{-3})$.
For the time test, replacing $\sqrt{1+r^2}$ by $N$ also produces a
vanishing error, because their difference and its first radial
derivative are $O(r^{-2})$.  The sphere area forms are the same.

The $g_0$ divergence of
$\psi\grad W_\mu-W_\mu\grad\psi$ is
$\psi LW_\mu-W_\mu L\psi$.  Apply the divergence theorem to
$[0,T]\times S^2$.  The outward normal of this integration domain
at its inner boundary is $-\partial_s$, yielding exactly the sign of
the boundary term in \eqref{ads:ana:green}.  Finally
$L\psi=O(r^{-2\tau})$ by \eqref{ads:ana:actual-equations}; hence the
bulk integrals converge, and $T\to\infty$ proves the identity.

For the asserted boundedness, $W_\mu=O(e^s)$ and
$\dd V_{g_0}=O(e^{2s})\,\dd s\,\dd\omega$.
The $W_\mu F$ term is dominated by a constant times
\[
 \|F\|_{2\beta}e^{(3-2\beta)s},
\]
which is integrable by \eqref{ads:ana:beta}.  The $uLW_i$ term is
dominated by $C\|u\|_\beta e^{-\beta s}$, and $LW_0=0$.
The boundary terms are bounded by the stated norms.
For finite coefficient jets, apply the same finite-radius identity
coefficientwise.  Lemma~\ref{ads:ana:coefficients} supplies their
differentiated asymptotics, and all nonlinear coefficient products
discarded from the flux are $O(r^{-6})$ with first derivatives.
Thus their coefficient flux is also
\eqref{ads:ana:green-functional}.
\end{proof}

\subsection{Approximation on the full space of seeds}

Let $\mathscr T_{\tau,\alpha}$ denote the real vector space of smooth
TT tensors in $C^{1,\alpha}_\tau$.  No boundary sign is imposed on
this space.  We shall construct coefficient functionals on it even
when no actual branch is being considered.

\begin{lemma}[Rotation approximation and coefficient continuity]
\label{ads:lem:rotation}
Every $q\in\mathscr T_{\tau,\alpha}$ is a limit, in $\|\cdot\|_\beta$,
of finite angular type tensors $q_\nu\in\mathscr T_{\tau,\alpha}$.
The approximants may be chosen as averages of rotated copies of $q$.
Their weighted norms are uniformly bounded, and the construction
preserves a nonnegative boundary value $q_{ss}|_S$ when that condition
holds.

The maps $q\mapsto u_1[q]$ and $q\mapsto u_2[q]$ defined for finite
angular type seeds extend uniquely by these approximations to
continuous, respectively linear and quadratic, maps on
$\mathscr T_{\tau,\alpha}$ with values in the weighted uniform norm.
The corresponding pairs $(F_j,B_j)$ in \eqref{ads:ana:first-two}
extend continuously in the norms $\|\cdot\|_{2\beta}$ and
$\|\cdot\|_{C^0(S)}$.
\end{lemma}

\begin{proof}
Let $U(R)q$ denote the pullback by a rotation $R\in SO(3)$, and
let $\dd R$ be normalized invariant volume.  Define
\[
 k_\nu(R)=c_\nu\left(\frac{1+\tr R}{4}\right)^\nu,
 \qquad \int_{SO(3)}k_\nu\,\dd R=1,
 \qquad q_\nu=\int_{SO(3)}k_\nu(R)U(R)q\,\dd R.
\]
The expression raised to the power $\nu$ lies in $[0,1]$ and attains
one only at the identity.  The normalized densities therefore
concentrate there.  To verify this directly, outside any fixed
neighborhood their base is at most $1-\delta$, whereas on a smaller
neighborhood of positive volume it is at least $1-\delta/2$;
the ratio of the corresponding integrals tends to zero.

The map $R\mapsto U(R)q$ is continuous in $\|\cdot\|_\beta$.
On a compact part of $M$ this follows from smoothness.  On the tail,
the bound $Ce^{-(\tau-\beta)s}$ is uniform in $R$, so the strict
weight margin controls it.  Concentration gives
$\|q_\nu-q\|_\beta\to0$.
Rotations preserve $g_0$, the end background, and $s$.  They commute
with trace and divergence and preserve the weighted regularity.
Consequently each average is smooth and TT, with the asserted uniform
bounds.  They also preserve the radial normal, so positivity of
$q_{ss}|_S$ is retained by the nonnegative average.

For fixed $\nu$, the density is a polynomial in matrix entries.
After a change of integration variable, the rotation orbit of $q_\nu$
depends on translates of this polynomial.  These translates lie in a
finite-dimensional polynomial space; their coefficients multiply
finitely many fixed averaged tensors.  Thus $q_\nu$ has finite angular
type.

It remains to check that the coefficient limits are independent of
the approximation.  On a bounded set in $\|q\|_\beta$, comparison
applied to \eqref{ads:ana:first-two} gives, for finite type seeds $q,\hat q$,
\begin{align*}
 \|u_1[q]-u_1[\hat q]\|_\beta
     &\leq C\|q-\hat q\|_\beta,\\
 \|F_2[q]-F_2[\hat q]\|_{2\beta}
  +\|B_2[q]-B_2[\hat q]\|_{C^0(S)}
     &\leq C\|q-\hat q\|_\beta,\\
 \|u_2[q]-u_2[\hat q]\|_\beta
     &\leq C\|q-\hat q\|_\beta.
\end{align*}
Here the elementary product estimate is
\[
 \bigl\||q|^2-|\hat q|^2\bigr\|_{2\beta}
 \leq(\|q\|_\beta+\|\hat q\|_\beta)\|q-\hat q\|_\beta.
\]
The comparison constant has no angular-cutoff dependence.
These estimates prove existence, uniqueness, and continuity of the
limits in the complete weighted uniform space.  Formula
\eqref{ads:ana:first-two} preserves their linear and quadratic dependence.
No derivative convergence of the limiting coefficients is needed in
this construction or in the mass functional.
\end{proof}

\subsection{Actual expansions and mass coefficients}

\begin{proposition}[Expansions of the given branch]
\label{ads:prop:expansions}
For every prescribed seed and branch in the theorem, set $p=2$ and
take $u_j,F_j,B_j$ from Lemma~\ref{ads:lem:rotation}.  If the seed has the
form \eqref{ads:ana:kernel-seed}, we may instead take $p=4$ with the
finite coefficient construction.  In both cases,
\begin{align}
 \left\|\psi-\sum_{j=1}^p\eps^ju_j\right\|_\beta
     &=O(\eps^{p+1}),\notag\\
 \left\|L\psi-\sum_{j=1}^p\eps^jF_j\right\|_{2\beta}
     &=O(\eps^{p+1}),\notag\\
 \left\|\partial_s\psi|_S-\sum_{j=1}^p\eps^jB_j\right\|_{C^0(S)}
     &=O(\eps^{p+1}).
 \label{ads:ana:three-remainders}
\end{align}
The area and all four mass components have actual expansions to the
same order.  Their mass coefficients are
\begin{equation}
 a_{\mu j}[q]=\mathcal G_\mu(u_j,F_j,B_j),\qquad
 p_\mu(\eps)=p_\mu(0)+\sum_{j=1}^p\eps^j a_{\mu j}[q]
                      +O(\eps^{p+1}).
 \label{ads:ana:mass-coefficients}
\end{equation}

In particular, the actual deficit satisfies
\begin{equation}
 \mathfrak D_q(\eps)=c_2[q]\eps^2+O(\eps^3)
 \label{ads:ana:deficit-second}
\end{equation}
for every seed under consideration.  The coefficient $c_2$ is a
continuous quadratic form on $\mathscr T_{\tau,\alpha}$ in the
$\|\cdot\|_\beta$ topology, invariant under rotations.  For seeds
of the form \eqref{ads:ana:kernel-seed} there is additionally the actual
expansion
\begin{equation}
 \mathfrak D_q(\eps)
     =c_2[q]\eps^2+c_3[q]\eps^3+c_4[q]\eps^4+O(\eps^5).
 \label{ads:ana:deficit-fourth}
\end{equation}
The subsequent geometric arguments determine the signs and the
vanishing of these coefficients.
\end{proposition}

\begin{proof}
\emph{Initial estimate.}
The given convergence of $\psi$ in $C^{2,\alpha}_\tau$ implies that
$\|\psi\|_{C^0}$ is small.  Write
$\mathcal A(\psi)=V_\psi\psi$, extending the quotient by zero at
$\psi=0$.  Then $\|V_\psi\|_{C^0}=o(1)$.  Moving this term to the
potential in \eqref{ads:ana:actual-equations} leaves an interior source
bounded by $C\eps^2e^{-2\tau s}$ and boundary derivative bounded
by $C\eps$.  Since $\psi=o(e^{-\beta s})$, Lemma~\ref{ads:lem:inverse}
gives
\begin{equation}
 \|\psi\|_\beta=O(\eps).
 \label{ads:ana:initial-bound}
\end{equation}

\emph{Finite angular type seeds.}
Suppose first that the needed coefficients have been constructed
by \eqref{ads:ana:coefficient-recursion}.  The first bound
\eqref{ads:ana:initial-bound} starts an induction.  If
\[
 \|\psi-P_{j-1}\|_\beta=O(\eps^j),
\]
then $\psi$ and $P_{j-1}$ are both $O(\eps e^{-\beta s})$.
For such arguments, the mean value theorem applied to
\eqref{ads:ana:nonlinearities} gives
\[
 |\mathcal N_{\eps,q}(\psi)
       -\mathcal N_{\eps,q}(P_{j-1})|
 \leq C\eps e^{-\beta s}|\psi-P_{j-1}|
 =O(\eps^{j+1}e^{-2\beta s}).
\]
Taylor substitution into the finite polynomial $P_{j-1}$ leaves
an error of this same order after its coefficients through degree $j$
are extracted.  Every interior term contains at least two decaying
factors, counting the factor $Q_q$ as two.  On $S$,
\[
 |\mathcal B_{\eps,q}(\psi)
       -\mathcal B_{\eps,q}(P_{j-1})|
 \leq C\eps|\psi-P_{j-1}|=O(\eps^{j+1}),
\]
and the boundary Taylor error has the same order.  Thus
$\psi-P_j$ has interior source
$O(\eps^{j+1}e^{-2\beta s})$ and boundary derivative
$O(\eps^{j+1})$.  It is $o(e^{-\beta s})$ at infinity, so comparison
improves its weighted norm to $O(\eps^{j+1})$.
This proves all three assertions in \eqref{ads:ana:three-remainders}
inductively through the required degree.  In particular, it proves
the fourth-order assertions in the special case
\eqref{ads:ana:kernel-seed}.

\emph{Arbitrary seeds through degree two.}
Let $q_\nu$ be the approximants of Lemma~\ref{ads:lem:rotation}, and
put $d_\nu=\|q-q_\nu\|_\beta$.  Write $u_{j,\nu},F_{j,\nu},B_{j,\nu}$
for their coefficients.  These have uniform bounds in the norms of
that lemma.  We compare them directly with the given branch for $q$;
no branch for $q_\nu$ is introduced.
Using \eqref{ads:ana:initial-bound}, the equation for
$\psi-\eps u_{1,\nu}$ has source
$O(\eps^2e^{-2\beta s})$ and boundary derivative
$O(\eps d_\nu+\eps^2)$.  Hence
\begin{equation}
 \|\psi-\eps u_{1,\nu}\|_\beta
       \leq C(\eps d_\nu+\eps^2).
 \label{ads:ana:approx-first}
\end{equation}
Expanding the nonlinearities once more, formula
\eqref{ads:ana:first-two} and \eqref{ads:ana:approx-first} imply
\begin{align}
 \|L\psi-\eps^2F_{2,\nu}\|_{2\beta}
   &\leq C(\eps^2d_\nu+\eps^3),\notag\\
 \|\partial_s\psi|_S-\eps B_{1,\nu}-\eps^2B_{2,\nu}\|_{C^0(S)}
   &\leq C\bigl((\eps+\eps^2)d_\nu+\eps^3\bigr).
 \label{ads:ana:approx-sources}
\end{align}
For clarity, the only quadratic interior discrepancies are
$\psi^2-\eps^2u_{1,\nu}^2$ and
$\eps^2(Q_q-Q_{q_\nu})$; their weighted norms are bounded by the
first right-hand side.  At the boundary the linear discrepancy is
$\eps(h_q-h_{q_\nu})/4$, and the next one contains
$\eps h_q\psi-\eps^2h_{q_\nu}u_{1,\nu}$, producing exactly the
second bound.  All remaining terms are cubic with the stated weights.
Comparison applied to
$\psi-\eps u_{1,\nu}-\eps^2u_{2,\nu}$ now gives
\[
 \|\psi-\eps u_{1,\nu}-\eps^2u_{2,\nu}\|_\beta
 \leq C\bigl((\eps+\eps^2)d_\nu+\eps^3\bigr).
\]
The constants are independent of $\nu$.  Letting $\nu\to\infty$
and using coefficient continuity proves the full three estimates
\eqref{ads:ana:three-remainders} for $p=2$ and the original arbitrary seed.

\emph{Mass, area, and the deficit.}
Apply the bounded functional in Lemma~\ref{ads:lem:green} to the three
remainders.  This proves \eqref{ads:ana:mass-coefficients} directly.
The area is the compact-boundary integral
\[
 A_\eps=\int_S(1+\psi)^4\,\dd A_{g_0},
\]
so its expansion follows from the first remainder alone.  If
$A_\eps=A_0+\eps A_1+\eps^2 A_2+O(\eps^3)$, then
\[
 A_1=4\int_Su_1\,\dd A_{g_0},\qquad
 A_2=\int_S(4u_2+6u_1^2)\,\dd A_{g_0}.
\]
The mass norm is smooth near $(m_0,0,0,0)$ and $b_*$ is smooth near
$A_0=4\pi a^2$.  Their expansions therefore have the same controlled
remainders.  The constant deficit is zero.  Also $F_1=0$,
$N|_S=0$, $\partial_sN|_S=\kappa$, and $LN=0$, so
\[
 a_{01}=\frac\kappa{2\pi}\int_Su_1\,\dd A_{g_0}
        =b_*'(A_0)A_1.
\]
The spatial mass components have zero base value and first enter
the mass norm quadratically.  This proves that the linear deficit
coefficient vanishes.

In degree two the scalar expression is explicitly
\begin{equation}
 c_2[q]=a_{02}
       -\frac1{2m_0}\sum_{i=1}^3a_{i1}^2
       -b_*'(A_0)A_2-\frac12b_*''(A_0)A_1^2.
 \label{ads:ana:quadratic-functional}
\end{equation}
Lemma~\ref{ads:lem:rotation} and the bounded Green functional show that
this is a continuous quadratic form on the entire formal seed space
$\mathscr T_{\tau,\alpha}$.  Formula \eqref{ads:ana:quadratic-functional}
uses only the coefficient functions and absolutely convergent Green
integrals; it does not assume a geometric integral formula for
arbitrary seeds.  Rotations leave the area invariant, fix the time
mass component, and rotate the three spatial components.  Thus $c_2$
is rotation invariant.  The identical finite coefficient argument
through degree four proves \eqref{ads:ana:deficit-fourth}.
\end{proof}

\begin{remark}
\label{ads:ana:order-of-limits}
The derivative control used in \eqref{ads:ana:flux-errors} comes from each
actual member's assumed $C^{2,\alpha}_\tau$ decay.  It is not a
consequence of the weighted uniform Taylor estimates.  The radius
limit is taken for each fixed member to obtain the exact identity
\eqref{ads:ana:green}; only then are its integrals expanded.  Similarly,
each finite angular type coefficient is identified with its own flux
before passing to an approximant limit in the bounded functional
\eqref{ads:ana:green-functional}.  No uniform convergence of discarded
derivative flux errors in either parameter or angular cutoff is used.
\end{remark}
